\section{Implementación}

\subsection{Arquitectura}

La app sigue MVVM con principios de Clean Architecture. Las dependencias apuntan hacia el dominio: la
interfaz conoce a los ViewModels, los ViewModels a los casos de uso y los casos de uso solo a interfaces
de repositorio. Room es la única fuente de verdad. La red actualiza Room y la interfaz observa Room
mediante \texttt{Flow}.

\begin{tabularx}{\linewidth}{@{}>{\ttfamily}l L@{}}
\toprule
Paquete & Responsabilidad \\
\midrule
ui & Pantallas Compose, componentes, tema y navegación \\
presentation & ViewModels y estados de pantalla \\
domain & Modelos, reglas de puntos, casos de uso e interfaces, en Kotlin puro \\
data & Room, DataStore, Retrofit, repositorios y mapeos \\
platform & Workers, UsageStatsManager, notificaciones, ubicación y tile \\
di & Contenedor de dependencias \\
\bottomrule
\end{tabularx}

El backend es una API en Ktor con Exposed y PostgreSQL. Tiene autenticación con JWT, el catálogo y el
libro de puntos, y la confirmación de canjes con PIN.

\subsection{Tecnologías}

Kotlin, Android, Jetpack Compose, Navigation Compose, ViewModel con Flow, Room, DataStore, Retrofit,
WorkManager, Ktor, Exposed y PostgreSQL. Las pruebas usan JUnit, Turbine, MockWebServer y work-testing.
Las pruebas del backend corren sobre Ktor \texttt{testApplication}.

\subsection{Decisiones relevantes}

\begin{itemize}
\item \textbf{Saldo confirmado por el servidor.} El teléfono informa el total de cada día y recibe el
saldo que se puede gastar. Lo ganado que todavía no se informó aparece aparte como ``por confirmar''
y no se puede gastar hasta que el servidor lo confirme. Si nunca hubo una sincronización, el saldo se
calcula en el teléfono.
\item \textbf{Idempotencia.} El teléfono genera el identificador del canje. Si la red falla y se reintenta,
el servidor devuelve el mismo canje y no lo descuenta dos veces. Lo mismo vale para el envío de puntos.
\item \textbf{Concurrencia en el servidor.} El descuento de saldo y de stock se hace dentro de la misma
transacción, con bloqueo sobre el usuario. Dos canjes simultáneos no pueden gastar el mismo saldo ni
tomar el último stock.
\item \textbf{Código de un solo uso.} El código vence a los 10 minutos. Un PIN incorrecto cuenta como
intento, y con 5 intentos fallidos el canje se rechaza y los puntos vuelven al saldo.
\item \textbf{Privacidad.} Ninguna pausa ni horario sale del teléfono. El servidor solo recibe totales
por día, con un máximo de 240 puntos. La app no registra nombres de otras aplicaciones, solo los
eventos de pantalla y de servicios que necesita.
\item \textbf{Tokens cifrados.} La sesión se guarda cifrada con AES-GCM y una clave del Android Keystore.
\item \textbf{Reglas configurables.} Los parámetros de puntos (horario de descanso, mínimos, topes,
bonus) viven en un objeto de configuración. Cambiarlos no obliga a tocar el código que los aplica.
\end{itemize}

\subsection{Sincronización offline}

El catálogo se sincroniza cada 6 horas con red, y al abrir Explorar si pasó más de una hora desde la
última sincronización. Los puntos se envían al abrir la app y después de cada validación que sumó. Si no
hay red, la sincronización espera a que vuelva. Los canjes y la confirmación con PIN necesitan conexión
y avisan al usuario cuando no la tienen.

\subsection{Cómo se ejecuta}

Las instrucciones completas están en el README de cada repositorio. En la práctica, la app se abre en
Android Studio, y el backend se levanta con \texttt{docker compose up -d --build}. La app de debug apunta
a \texttt{http://10.0.2.2:8080/}.
